%----------------------------------------------------------------------------------------
%	PACKAGES AND OTHER DOCUMENT CONFIGURATIONS
%----------------------------------------------------------------------------------------

\documentclass[10pt,stdletter,dateno,sigleft]{newlfm} % Extra options: 'sigleft' for a left-aligned signature, 'stdletternofrom' to remove the from address, 'letterpaper' for US letter paper - consult the newlfm class manual for more options

%\usepackage{charter} % Use the Charter font for the document text
\usepackage{graphics}

\usepackage{standalone}

\newsavebox{\Luiuc}\sbox{\Luiuc}{\parbox[b]{1.75in}{\vspace{0.5in}
\includegraphics[width=1.2\linewidth]{/Users/hectorbahamonde/RU/Job_Market/logo.jpg}}} % Company/institution logo at the top left of the page
\makeletterhead{Uiuc}{\Lheader{\usebox{\Luiuc}}}

\newlfmP{sigsize=50pt} % Slightly decrease the height of the signature field
\newlfmP{addrfromphone} % Print a phone number under the sender's address
\newlfmP{addrfromemail} % Print an email address under the sender's address
\PhrPhone{\texttt{p}} % Customize the "Telephone" text
\PhrEmail{\texttt{e}} % Customize the "E-mail" text

\lthUiuc % Print the company/institution logo

% email input begin
\newread\fid
\newcommand{\readfile}[1]% #1 = filename
{\bgroup
  \endlinechar=-1
  \openin\fid=#1
  \read\fid to\filetext
  \loop\ifx\empty\filetext\relax% skip over comments
    \read\fid to\filetext
  \repeat
  \closein\fid
  \global\let\filetext=\filetext
\egroup}
\readfile{/Users/hectorbahamonde/RU/Bibliografia_PoliSci/email.txt}
% email input end

\usepackage{color} % for colors
\usepackage{hyperref} % for links

\hypersetup{
  colorlinks = true,
  urlcolor = blue,
  pdfpagemode = UseNone
}

%----------------------------------------------------------------------------------------
%	YOUR NAME AND CONTACT INFORMATION
%----------------------------------------------------------------------------------------

%\namefrom{\vspace{-3cm}Hector Bahamonde, PhD} % Name

\addrfrom{
{\vspace{-2cm}}\\ % Date
[12pt]
{\bf Hector Bahamonde, PhD}\\
\color[HTML]{003300}{{\bf Postdoctoral Fellow}}\\
\emph{\color[HTML]{003300}{Center For Inter-American Policy \& Research}}\\
\color[HTML]{003300}{Tulane University}\\
205 Richardson, New Orleans, LA 70118
}

\phonefrom{(732) 318-9650} % Phone number

\emailfrom{\href{mailto:\filetext}{\filetext} \\ 
\texttt{w}: \href{http://www.hectorbahamonde.com}{www.HectorBahamonde.com}\\
\today, download latest version \href{http://github.com/hbahamonde/Job_Market/raw/master/Bahamonde_Cover_Letter.pdf}{\texttt{{\color{red}here}}}}% Link to last version


%----------------------------------------------------------------------------------------
%	ADDRESSEE AND GREETING/CLOSING
%----------------------------------------------------------------------------------------

%\greetto{\vspace{-1cm}Dear Members of the Search Committee,} % Greeting text
%\closeline{\vspace{-1cm}\includegraphics[width=3cm]{/Users/hectorbahamonde/RU/Administracion/signature.pdf} \hspace{-4cm}} % Closing text




%----------------------------------------------------------------------------------------

\begin{document}
\begin{newlfm}

%----------------------------------------------------------------------------------------
%	LETTER CONTENT
%----------------------------------------------------------------------------------------



\vspace{-2cm}{\bf \huge Diversity Statement}\\

Rutgers is one of the most diverse universities in the United States. Being a teaching assistant of large classes has given me the opportunity to experience diversity. Diverse environments not only make individuals more tolerant, flexible and educated. They also force us, in a good way, to develop teaching philosophies that are flexible, considerate and interesting for diverse student bodies. 

There are many types of diversity, such as economic, cultural, sexual and political. We, as social scientists, should know that all these types are part of our daily life as researchers and instructors. Everything we do, will most likely be related to economic inequality, racial politics or gender representation. Everything we say is subject to debate and tension. These disputes should be treated even more carefully, specially if our audience is diverse. It is not about not making political jokes. It is more about answering the following questions: \emph{How can I make today's lecture interesting enough, so my students engage not only with me but with all their classmates? How should I present this contentious issue related to race, for example, in a way that at the end of class, my students know that there is still much more to be done?} In my experience, the answers have less to do with teaching ``the facts.'' They have to do with knowing how to expose students to complex issues. \emph{Should the rich be taxed to help the poor? Was the atomic bomb a reasonable measure? Are all democratic outcomes ``good''?} (Hitler was democratically elected) \emph{Should rich states help poor states?} (the economic literature says ``no''), etc. All these topics can be answered from so many angles. Moreover, \emph{students will answer based on their own economic, cultural, sexual and political backgrounds}. At the end of the day, I know that I did my job well if I know that my students felt \emph{curious} about these issues, even at the cost of leaving the classroom in an atmosphere of ``uncertainty.''

All and all, as instructors, I think diversity plays always in our favor. By all means, it makes lecturing more interesting. But it also plays in favor of our  students. Diversity exposes them to different perspectives. However, even more importantly, diversity forms better citizens better able to engage in a diverse world. That is why I think we should always ``take advantage'' of diversity.

% ending
{More information, \href{http://www.hectorbahamonde.com/teaching/}{\texttt{syllabi}}, my \href{http://github.com/hbahamonde/Job_Market/raw/master/Bahamonde_Research_Statement.pdf}{\texttt{research}}, \href{http://github.com/hbahamonde/Job_Market/raw/master/Bahamonde_Teaching_Statement.pdf}{\texttt{teaching}} and \href{http://github.com/hbahamonde/Job_Market/raw/master/Bahamonde_Diversity_Statement.pdf}{\texttt{diversity}} statements, as well as other \href{http://www.hectorbahamonde.com/research/}{\texttt{papers}} are available on my website: \href{http://www.hectorbahamonde.com}{\texttt{www.HectorBahamonde.com}}. Thank you for considering my application. I look forward to hearing from you.\unskip}


\end{newlfm}
\end{document}

